
\begin{example}
    点 $P(-2,-1)$ 到直线 $l:(1+3 \lambda) x+(1+\lambda) y-2-4 \lambda=0(\lambda \in \mathrm{R})$ 的距离最大时，其最大值以及此时的直线方程分别为 ()
    A．$\sqrt{13} ; 3 x+2 y-5=0$
    B．$\sqrt{11} ; 3 x+2 y-5=0$
    C．$\sqrt{13} ; 2 x-3 y+1=0$
    D．$\sqrt{11} ; 2 x-3 y+1=0$

    2．直线 $y=k(x-5)-2(k \in \mathbf{R})$ 与圆 $(x-3)^2+(y+1)^2=6$ 的位置关系为 （\qquad）
    \begin{tasks}(4)
        \task 相离
        \task 相交
        \task 相切
        \task 无法确定
    \end{tasks}
\end{example}

\begin{example}
    已知圆 $C:(x-1)^2+(y-2)^2=25$ ，直线 $l:(2 m+1) x+(m+1) y-7 m-4=0$ ．则下列命题正确的有 （\qquad）
    \begin{tasks}(4)
        \task 直线 $l$ 恒过定点 $(3,1)$
        \task 圆 $C$ 被 $y$ 轴截得的弦长为 $2 \sqrt{6}$
        \task 直线 $l$ 与圆 $C$ 恒相交
        \task 直线， 1 被圆 $C$ 截得弦长最短时，直线 $l$ 的方程为 $2 x-y-5=0$
    \end{tasks}
\end{example}
\begin{example}
    已知直线 $l: y=k x-k$ ，圆 $C: x^2+y^2=4$ ，则下列结论正确的有 （\qquad）
    \begin{tasks}(4)
        \task 直线 $l$ 过定点 $(1,0)$
        \task 直线 $l$ 与圆 $C$ 恒相交
        \task 直线 $l$ 被圆 $C$ 截得的弦长最短为 4
        \task 若直线 $l$ 被圆 $C$ 截得的弦长为 $\sqrt{14}$ ，则 $k= \pm 1$
    \end{tasks}
\end{example}
三、填空题

5．若直线 $k x-y-2 k+3=0$ 必过一定点，则该定点坐标是 $\qquad$ ．

6．已知圆 $C:(x+2)^2+(y-4)^2=1$ ，则圆心 $C$ 到直线 $l: k x+y-k=0$ 的最大距离为 $\qquad$ ．

四、解答题

\begin{example}
    已知直线 $l: k x-y+2+k=0(k \in \mathbf{R})$ ．
    \begin{enumerate}
        \item 若直线不经过第三象限，求 $k$ 的取值范围；

        \item 若直线 $l$ 交 $x$ 轴负半轴于 $A$ ，交 $y$ 轴正半轴于 $B, \triangle A O B$ 的面积为 $S$（ $O$ 为坐标原点），求 $S$ 的最小值和此时直线 $l$ 的方程．
    \end{enumerate}
\end{example}

\begin{example}
    已知直线 $l_1: x+m y+6=0, ~ l_2:(m-2) x+3 y+2 m=0$ ，根据下列条件分别求实数 $m$ 的取值范围．
    \begin{enumerate}
        \item $l_1$ 与 $l_2$ 相交；

        \item $l_1$ 与 $l_2$ 重合．
    \end{enumerate}
\end{example}

\begin{example}
    已知两条直线 $l_1: x+m^2 y+6=0, l_2:(m-2) x+3 m y+2 m=0$ ，
    \begin{enumerate}
        \item 当 $m$ 为何值时，$l_1$ 与 $l_2$ 相交；

        \item $l_1$ 与 $l_2$ 是两条不同直线，$l_2$ 经过定点 $A$ ，当 $l_1$ 也经过点 $A$ 时，求 $m$ 的值．
    \end{enumerate}
\end{example}

\begin{example}
    已知直线 $l:(2 m+1) x+(m+1) y=8 m+5$ ，圆 $C:(x-1)^2+(y-2)^2=25$ ．
    \begin{enumerate}
        \item 证明：直线与圆总有两个交点，与 $m$ 的取值无关．

        \item 是否存在 $m$ ，使得直线 $/$ 被圆 $C$ 截得的弦长为 $2 \sqrt{22}$ ，若存在，求出 $m$ 的值；若不存在，请说明理由．
    \end{enumerate}
\end{example}

\begin{example}
    已知圆 $C$ 过三点 $(1,3),(1,-7),(4,2)$ ．
    \begin{enumerate}
        \item 求圆 $C$ 的方程；

        \item 斜率为 1 的直线 $l$ 与圆 $C$ 交于 $M, N$ 两点，若 $\triangle C M N$ 为等腰直角三角形，求直线 $l$ 的方程．
    \end{enumerate}
\end{example}